\documentclass[10pt,a4paper]{amsart}
\usepackage{amsmath,amssymb,amsthm,graphicx}
% For commutative diagrams we use Xy-pic
\usepackage[all]{xy}

% We use 2cell for 2-commutative diagrams.
\xyoption{2cell}
% Theorem environments.
%
\theoremstyle{plain}
\newtheorem{theorem}[subsection]{Theorem}
\newtheorem{proposition}[subsection]{Proposition}
\newtheorem{lemma}[subsection]{Lemma}

\theoremstyle{definition}
\newtheorem{definition}[subsection]{Definition}
\newtheorem{example}[subsection]{Example}
\newtheorem{exercise}[subsection]{Exercise}
\newtheorem{situation}[subsection]{Situation}

\theoremstyle{remark}
\newtheorem{remark}[subsection]{Remark}
\newtheorem{remarks}[subsection]{Remarks}

\numberwithin{equation}{subsection}

% Macros
%
\def\lim{\mathop{\mathrm{lim}}\nolimits}
\def\colim{\mathop{\mathrm{colim}}\nolimits}
\def\Spec{\mathop{\mathrm{Spec}}}
\def\Hom{\mathop{\mathrm{Hom}}\nolimits}
\def\Ext{\mathop{\mathrm{Ext}}\nolimits}
\def\SheafHom{\mathop{\mathcal{H}\!\mathit{om}}\nolimits}
\def\SheafExt{\mathop{\mathcal{E}\!\mathit{xt}}\nolimits}
\def\Sch{\mathit{Sch}}
\def\Mor{\mathop{\mathrm{Mor}}\nolimits}
\def\Ob{\mathop{\mathrm{Ob}}\nolimits}
\def\Sh{\mathop{\mathit{Sh}}\nolimits}
\def\NL{\mathop{N\!L}\nolimits}
\def\CH{\mathop{\mathrm{CH}}\nolimits}
\def\proetale{{pro\text{-}\acute{e}tale}}
\def\etale{{\acute{e}tale}}
\def\QCoh{\mathit{QCoh}}
\def\Ker{\mathop{\mathrm{Ker}}}
\def\Im{\mathop{\mathrm{Im}}}
\def\Coker{\mathop{\mathrm{Coker}}}
\def\Coim{\mathop{\mathrm{Coim}}}

% Boxtimes
%
\DeclareMathSymbol{\boxtimes}{\mathbin}{AMSa}{"02}

%
% Macros for moduli stacks/spaces
%
\def\QCohstack{\mathcal{QC}\!\mathit{oh}}
\def\Cohstack{\mathcal{C}\!\mathit{oh}}
\def\Spacesstack{\mathcal{S}\!\mathit{paces}}
\def\Quotfunctor{\mathrm{Quot}}
\def\Hilbfunctor{\mathrm{Hilb}}
\def\Curvesstack{\mathcal{C}\!\mathit{urves}}
\def\Polarizedstack{\mathcal{P}\!\mathit{olarized}}
\def\Complexesstack{\mathcal{C}\!\mathit{omplexes}}
% \Pic is the operator that assigns to X its picard group, usage \Pic(X)
% \Picardstack_{X/B} denotes the Picard stack of X over B
% \Picardfunctor_{X/B} denotes the Picard functor of X over B
\def\Pic{\mathop{\mathrm{Pic}}\nolimits}
\def\Picardstack{\mathcal{P}\!\mathit{ic}}
\def\Picardfunctor{\mathrm{Pic}}
\def\Deformationcategory{\mathcal{D}\!\mathit{ef}}
\newenvironment{reference}{}{}
\newenvironment{slogan}{}{}
\newenvironment{history}{}{}
\usepackage{cleveref}
\crefname{originallemma}{original lemma}{original lemmas}
\Crefname{originallemma}{Original lemma}{Original lemmas}
\makeatletter
\newlabel{lemma-divisor-special-fiber}{{0C5Z}{source}{}{sourcepointer-1}{}}
\newlabel{lemma-divisor-special-fiber@cref}{{[originallemma][0][]0C5Z}{[1][0][]source}{}{}{}}
\makeatother
\crefname{originaldefinition}{original definition}{original definitions}
\Crefname{originaldefinition}{Original definition}{Original definitions}
\makeatletter
\newlabel{definition-picard-group}{{0C7H}{source}{}{sourcepointer-2}{}}
\newlabel{definition-picard-group@cref}{{[originaldefinition][0][]0C7H}{[1][0][]source}{}{}{}}
\makeatother
\crefname{originallemma}{original lemma}{original lemmas}
\Crefname{originallemma}{Original lemma}{Original lemmas}
\makeatletter
\newlabel{lemma-regular-model-pic}{{0C63}{source}{}{sourcepointer-3}{}}
\newlabel{lemma-regular-model-pic@cref}{{[originallemma][0][]0C63}{[1][0][]source}{}{}{}}
\makeatother
\crefname{originallemma}{original lemma}{original lemmas}
\Crefname{originallemma}{Original lemma}{Original lemmas}
\makeatletter
\newlabel{varieties-lemma-divisible}{{0C6P}{source}{}{sourcepointer-4}{}}
\newlabel{varieties-lemma-divisible@cref}{{[originallemma][0][]0C6P}{[1][0][]source}{}{}{}}
\makeatother
\title{The canonical map from the Picard group of the curve to the Picard group of the numerical type, and its kernel}
\author{The Stacks Project Authors}
\date{stacks-project:models.tex:local-snapshot-20261012; GNU FDL 1.2 or later}
\begin{document}
\maketitle
\noindent\textit{Editorial note.} The counted unit is the Stacks Project Lemma lemma-canonical-map-of-pic at source lines 5911-5923 together with its proof at source lines 5925-5960, a straight adjacent pairing in the chapter file. The chapter is the upstream file models.tex, which is the file name recorded in the source record; the pool directory carries the longer name stacks-models and is not used as the chapter name anywhere in this card.\par\medskip
\noindent\textit{Editorial note.} The retained statement opens with the words In Situation followed by a reference to situation-regular-model, so the situation environment that defines the setting is retained as a context span at source lines 4628-4637. It is context, not a counted unit, and it does not overlap the counted interval.\par\medskip
\noindent\textit{Editorial note.} This chapter's candidate index records the statement body and the proof body without their closing markers, so the retained spans used here are the full environments, exactly as in the earlier cards of this chapter.\par\medskip
\noindent\textit{Editorial note.} The retained proof draws a commutative diagram of two exact rows using the author's Xy-pic matrix syntax, so the retained preamble keeps the chapter preamble's Xy-pic block at lines 14-18 verbatim, including the option that enables two-cell diagrams.\par\medskip
\noindent\textit{Editorial note.} One asymmetry of the author's file is retained verbatim rather than repaired: in the first line of the proof the author writes the closing parenthesis of the structure sheaf of the component inside the braces of the index, which is the author's own typography. Nothing has been normalised.\par\medskip
\noindent\textit{Editorial note.} The retained proof quotes one definition of the same chapter, one lemma of the same chapter, and one lemma of the Varieties chapter. All are declared as navigation pointers so that the references resolve; every one of them carries a stable Stacks tag in the shared tag index and those tags are recorded as the numbers, so the lane's fallback of recording the label string itself was not needed for this card.\par\medskip
\noindent\textit{Editorial note.} No citation occurs in the retained spans, so no bibliography entry is supplied, and the retained preamble is the chapter preamble's Xy-pic setup at lines 14-18, the theorem declarations and mathematical macros at lines 154-217, and the three no-op wrapper environments.\par\medskip
\noindent\textit{Allowed standard prerequisites (editorial).}\begin{itemize}
\item The numerical type of a model, divisibility of the degree by the residue field degree, and the exact sequence of the regular model. For a regular model the numerical type records the multiplicities, the genera and the intersection numbers of the components; the degree of any invertible sheaf on a component is divisible by the degree of the residue field extension of that component, which is what makes the divided-degree map well defined; and the Picard group of the model fits into the exact sequence determined by the multiplicities with the Picard group of the curve as its quotient. The author's own definitions and the two lemmas of this chapter that the proof quotes by name are the allowed inputs.. The counted map is the bridge used in this section to compare the Picard group of the curve with the combinatorial Picard group of the numerical type, and its kernel description is what turns the comparison into an exact sequence of torsion groups.
\end{itemize}\medskip
\ifdefined\hypertarget\hypertarget{sourcepointer-1}{}\fi
\noindent\textit{Source reference.} Original lemma 0C5Z: focused-seeds/AGP/sources/stacks-snapshot-20261011/tags.txt, source lines 15279--15279. Navigation pointer inside the retained situation, which names it when it writes the special fibre as a sum of its components with multiplicities. The number is the stable Stacks tag of the target, looked up in the shared tag index.\par\medskip
\ifdefined\hypertarget\hypertarget{sourcepointer-2}{}\fi
\noindent\textit{Source reference.} Original definition 0C7H: focused-seeds/AGP/sources/stacks-snapshot-20261011/tags.txt, source lines 15332--15332. Pointer to the author's own definition of the Picard group of a numerical type, quoted inside the retained proof; it is an allowed input of the section and is not reproduced inside the span. The number is the stable Stacks tag of the target, looked up in the shared tag index.\par\medskip
\ifdefined\hypertarget\hypertarget{sourcepointer-3}{}\fi
\noindent\textit{Source reference.} Original lemma 0C63: focused-seeds/AGP/sources/stacks-snapshot-20261011/tags.txt, source lines 15283--15283. Pointer to the exact sequence of the regular model, invoked by the retained proof for the exactness of the upper row of the diagram; that exact sequence is counted by a separate card of this lane and is not reproduced here. The number is the stable Stacks tag of the target (0C63), looked up in the shared tag index.\par\medskip
\ifdefined\hypertarget\hypertarget{sourcepointer-4}{}\fi
\noindent\textit{Source reference.} Original lemma 0C6P: focused-seeds/AGP/sources/stacks-snapshot-20261011/tags.txt, source lines 15304--15304. Allowed input of the Varieties chapter quoted by the retained proof for the divisibility of the degree of an invertible sheaf on a component by its residue field degree; the pointer is registered so the reference resolves and the cited statement is not reproduced inside the span. The number is the stable Stacks tag of the target, looked up in the shared tag index.\par\medskip
\begin{situation}
\label{situation-regular-model}
Let $R$ be a discrete valuation ring with fraction field $K$,
residue field $k$, and uniformizer $\pi$.
Let $C$ be a smooth projective curve over $K$ with $H^0(C, \mathcal{O}_C) = K$.
Let $X$ be a regular proper model of $C$.
Let $C_1, \ldots, C_n$ be the irreducible components of the special
fibre $X_k$. Write $X_k = \sum m_i C_i$ as in
Lemma \ref{lemma-divisor-special-fiber}.
\end{situation}

\begin{lemma}
\label{lemma-canonical-map-of-pic}
In Situation \ref{situation-regular-model}
let $T$ be the numerical type associated to $X$.
There exists a canonical map
$$
\Pic(C) \to \Pic(T)
$$
whose kernel is exactly those invertible modules on $C$
which are the restriction of invertible modules $\mathcal{L}$
on $X$ with $\deg_{C_i}(\mathcal{L}|_{C_i}) = 0$ for
$i = 1, \ldots, n$.
\end{lemma}

\begin{proof}
Recall that $w_i = [\kappa_i : k]$ where
$\kappa_i = H^0(C_i, \mathcal{O}_{C_i)})$ and recall
that the degree of any invertible module on $C_i$ is divisible
by $w_i$ (Varieties, Lemma \ref{varieties-lemma-divisible}).
Thus we can consider the map
$$
\frac{\deg}{w} : \Pic(X) \to \mathbf{Z}^{\oplus n}, \quad
\mathcal{L} \mapsto
(\frac{\deg(\mathcal{L}|_{C_1})}{w_1}, \ldots,
\frac{\deg(\mathcal{L}|_{C_n})}{w_n})
$$
The image of $\mathcal{O}_X(C_j)$ under this map is
$$
((C_j \cdot C_1)/w_1, \ldots, (C_j \cdot C_n)/w_n) =
(a_{1j}/w_1, \ldots, a_{nj}/w_n)
$$
which is exactly the image of the $j$th basis vector under the map
$(a_{ij}/w_i) : \mathbf{Z}^{\oplus n} \to \mathbf{Z}^{\oplus n}$
defining the Picard group of $T$, see
Definition \ref{definition-picard-group}.
Thus the canonical map of the lemma comes from the commutative
diagram
$$
\xymatrix{
\mathbf{Z}^{\oplus n} \ar[r] \ar[d]_{\text{id}} &
\Pic(X) \ar[r] \ar[d]^{\frac{\deg}{w}} &
\Pic(C) \ar[r] \ar[d] & 0 \\
\mathbf{Z}^{\oplus n} \ar[r]^{(a_{ij}/w_i)} &
\mathbf{Z}^{\oplus n} \ar[r] &
\Pic(T) \ar[r] & 0
}
$$
with exact rows (top row by Lemma \ref{lemma-regular-model-pic}).
The description of the kernel is clear.
\end{proof}

\end{document}
